\chapter*{Glosario}
\addcontentsline{toc}{chapter}{Glosario}

Los términos de este glosario se marcan con \textsuperscript{\textdagger} en su primera aparición en el texto, con un enlace a su definición.

% Mantener en orden alfabético. Cada término se marca con \glos{id}{texto} en su primera
% aparición en los capítulos; el id es el primer argumento de \termino.
\begin{description}
    \termino{federado}{Aprendizaje federado} Técnica de aprendizaje automático en la que varios dispositivos entrenan un modelo compartido a partir de sus datos locales y envían solo las actualizaciones del modelo, sin transferir los datos crudos.
    \termino{cabeza}{Cabeza de clasificación} Parte final del modelo, formada por las capas densas, que transforma el resumen de la secuencia en la probabilidad de cada clase. Es la parte que se entrena en el dispositivo.
    \termino{calibracion}{Calibración} Ajuste de las probabilidades que entrega un modelo para que la confianza informada se corresponda con su frecuencia real de acierto.
    \termino{contrato}{Contrato de características} Especificación única, compartida entre el entrenamiento y la aplicación, de cómo se transforma una toma en la entrada del modelo, con casos de prueba que ambas implementaciones deben reproducir.
    \termino{cuantizacion}{Cuantización} Técnica que reduce la precisión numérica de los parámetros de un modelo (por ejemplo, de punto flotante de 32 bits a enteros de 8 bits) para disminuir su tamaño y acelerar la inferencia.
    \termino{glosa}{Glosa} Transcripción de una seña mediante una palabra de una lengua oral, utilizada como representación intermedia en la traducción de lenguas de señas.
    \termino{handedness}{Handedness} Lateralidad de una mano detectada, es decir, si se trata de la mano izquierda o de la derecha.
    \termino{inferencia}{Inferencia} Uso de un modelo ya entrenado para obtener una predicción a partir de una entrada nueva.
    \termino{landmark}{Landmark} Punto de referencia que representa la posición de una articulación o extremo de la mano dentro de la imagen.
    \termino{pipeline}{Pipeline} Secuencia de etapas de procesamiento en la que la salida de cada etapa es la entrada de la siguiente.
    \termino{aislada}{Seña aislada} Seña individual que se reconoce de forma independiente, sin considerar las señas anteriores o posteriores dentro de una oración.
    \termino{toma}{Toma} Grabación de una única realización de una seña, desde que aparecen las manos hasta que se completa la captura.
    \termino{lopo}{Validación por persona} Esquema de evaluación en el que el modelo se entrena sin las tomas de una persona y se evalúa con ellas, para medir su desempeño con personas que no vio.
\end{description}
